\documentclass[journal=acsodf,manuscript=suppinfo]{achemso}
\usepackage{amsmath,booktabs,graphicx,siunitx}
\author{Krishna Harish}
\affiliation{Lawrence E. Elkins High School, Missouri City, Texas, USA}
\email{krishnaharish2009@gmail.com}
\title{Supporting Information: Objective-Dependent Active Learning and Calibrated
Uncertainty for Sample-Efficient Discovery of Rare-Earth Extractants}
\begin{document}

\section{Datasets}
\textbf{$f$-element extraction (primary).} 1{,}202 experimental lanthanide
distribution coefficients ($\log D$) from Liu et al.\ (1{,}085 training + 117
validation), with 2{,}291 features per record: ligand ECFP/Morgan bits and RDKit
physicochemical descriptors, experimental conditions (ligand and acid
concentration, diluent identity and volume ratio, temperature), and
lanthanide-identity descriptors. Target range $\log D\in[-4.00,3.59]$. A
strictly prospective external set of four subsequently synthesised ligands
(56 ligand$\times$metal measurements) is used only for final validation.

\textbf{Lipophilicity (generality).} 4{,}200 octanol--water $\log D$ values
(MoleculeNet). Features: 1{,}024-bit Morgan fingerprints (radius 2) concatenated
with the RDKit physicochemical descriptors above.

\section{Methods}
\textbf{Surrogate.} Random forest, 300 trees, minimum 2 samples per leaf, features
standardised. Epistemic uncertainty $\hat\sigma$ is the per-point standard deviation
across trees.

\textbf{Conformal calibration.} Split-conformal: 20\% of the training data held out
as a calibration set; nonconformity scores $s_i=|y_i-\hat\mu_i|/\hat\sigma_i$; the
calibrated interval half-width is $q_\alpha\hat\sigma$ with $q_\alpha$ the empirical
$\alpha$-quantile of $\{s_i\}$. Within the active-learning loop the calibration set is
drawn from the already-revealed pool at each budget and is disjoint from the fixed
held-out test set, so the rescaling never uses test data and never leaks from the
unrevealed pool. The subgroup-coverage analysis (main text, Table~3) additionally uses
a ligand-disjoint train/calibration/test partition, so calibration, training, and
testing share no ligand.

\textbf{Active-learning protocol.} Fixed $80/20$ pool/test split. $f$-element:
seed $=25$, batch $=15$, budget $=420$. Lipophilicity: seed $=40$, batch $=25$,
budget $=715$. Six random seeds per strategy; curves report mean$\pm$s.d. Top-10
recall is the fraction of the ten highest-$\log D$ pool members acquired.

\section{Full results}

\begin{table}[h]\centering
\caption{Active learning on the $f$-element $\log D$ benchmark (pool $962$,
test $240$). $R^2$ at fixed budget; final $R^2$ and top-10 recall at budget $430$;
measurements to reach $R^2\!=\!0.85$.}
\begin{tabular}{lccccc}
\toprule
strategy & $R^2$@100 & $R^2$@250 & final $R^2$ & recall & labels$\to$0.85 \\
\midrule
random      & 0.543 & 0.754 & 0.851 & 0.50 & 404 \\
greedy      & 0.272 & 0.506 & 0.633 & 1.00 & n/a \\
uncertainty & 0.396 & 0.679 & 0.870 & 0.38 & 400 \\
UCB         & 0.438 & 0.600 & 0.701 & 1.00 & n/a \\
EI          & 0.379 & 0.664 & 0.776 & 0.98 & n/a \\
hybrid      & 0.379 & n/a    & 0.782 & 0.93 & n/a \\
\bottomrule
\end{tabular}
\end{table}

\begin{table}[h]\centering
\caption{Active learning on the Lipophilicity benchmark (pool $3360$, test $840$),
Morgan-fingerprint features. Final values at budget $715$.}
\begin{tabular}{lcccc}
\toprule
strategy & $R^2$@150 & final $R^2$ & recall \\
\midrule
random      & 0.253 & 0.497 & 0.27 \\
greedy      & 0.105 & 0.206 & 0.53 \\
uncertainty & 0.212 & 0.501 & 0.20 \\
UCB         & 0.139 & 0.281 & 0.57 \\
EI          & 0.143 & 0.309 & 0.53 \\
hybrid      & 0.165 & 0.384 & 0.30 \\
\bottomrule
\end{tabular}
\end{table}

\begin{table}[h]\centering
\caption{Uncertainty calibration ($f$-element benchmark, 10 seeds). Empirical
coverage of nominal-$\alpha$ intervals.}
\begin{tabular}{lcc}
\toprule
nominal $\alpha$ & uncalibrated & conformal \\
\midrule
0.50 & 0.656 & 0.503 \\
0.80 & 0.904 & 0.778 \\
0.90 & 0.950 & 0.886 \\
\midrule
mean calib.\ error & 0.103 & 0.013 \\
\bottomrule
\end{tabular}
\end{table}

\begin{table}[h]\centering
\caption{Active learning with vs.\ without conformal calibration (final $R^2$ /
top-10 recall). Acquisition performance is unchanged within noise.}
\begin{tabular}{lcc}
\toprule
strategy & uncalibrated & conformal \\
\midrule
UCB    & 0.647 / 1.00 & 0.635 / 1.00 \\
EI     & 0.696 / 0.98 & 0.709 / 0.98 \\
hybrid & 0.706 / 0.90 & 0.706 / 0.90 \\
\bottomrule
\end{tabular}
\end{table}

\section{Baseline and prospective validation}
A random forest trained on the published training split attains $R^2=0.942$,
RMSE$=0.328$ on the validation split (cf.\ the reported deep network $R^2=0.85$,
RMSE$=0.53$). Trained on all 1{,}202 measurements and applied to the four
prospective ligands, it attains $R^2=0.687$, RMSE$=0.639$, MAE$=0.502$ log units.

\section{Revision experiments}
The following tables collect the full results of the leakage-controlled evaluation,
the definition-resolved discovery metrics, the subgroup-resolved calibration, and the
model/uncertainty-estimator comparison added in revision. Every value is generated
from the released results files by \texttt{revision/make\_tables.py}.

\begin{table}[htbp]
\centering
\caption{Random-forest generalisation under increasingly stringent, leakage-controlled splits. Random row-level cross-validation is strongly optimistic; grouped splits that forbid a ligand, source, or structural family from appearing in both train and test collapse the estimate toward the mean. Mean$\pm$s.d.\ over 4 repeats of 5-fold grouped CV.}
\label{tab:splits}
\begin{tabular}{lccc}
\toprule
Split & $R^2$ & RMSE & MAE \\
\midrule
Random (row-level) & 0.88$\pm$0.03 & 0.46 & 0.28 \\
Ligand-disjoint & 0.20$\pm$0.31 & 1.14 & 0.87 \\
Source-disjoint & -0.03$\pm$0.80 & 1.22 & 0.98 \\
Ligand-family (scaffold proxy) & -0.24$\pm$0.50 & 1.31 & 1.04 \\
\midrule
Prospective (4 new ligands) & 0.70 & 0.63 & 0.49 \\
\bottomrule
\end{tabular}
\end{table}

\begin{table}[htbp]
\centering
\caption{Discovery performance at the final budget under three definitions of a `top-10 system' (condition-specific row, ligand--metal pair, unique ligand), plus simple regret (log units) and hit-rate above a practical threshold ($\log D\ge 1$, i.e.\ $D\ge 10$). The ranking of strategies inverts with the definition: exploitation dominates for rows, but broad sampling already recovers most top ligands. Means over six seeds.}
\label{tab:recall}
\begin{tabular}{lccccc}
\toprule
Strategy & recall & recall & recall & regret & hit-rate \\
 & (row) & (lig--metal) & (ligand) & (log) & ($\log D\!\ge\!1$) \\
\midrule
random & 0.45 & 0.72 & 0.95 & 0.15 & 0.26 \\
greedy & 1.00 & 0.88 & 0.82 & 0.00 & 0.60 \\
uncertainty & 0.42 & 0.85 & 0.93 & 0.12 & 0.26 \\
ucb & 1.00 & 0.90 & 0.88 & 0.00 & 0.59 \\
ei & 0.93 & 0.90 & 0.97 & 0.02 & 0.47 \\
hybrid & 0.80 & 0.87 & 0.97 & 0.01 & 0.52 \\
\bottomrule
\end{tabular}
\end{table}

\begin{table}[htbp]
\centering
\caption{Uncertainty calibration under a realistic ligand-disjoint split (nominal coverage 80\%). Raw ensemble intervals under-cover; a single global split-conformal factor restores average coverage but leaves subgroup imbalance, which group-conditional (Mondrian) conformal reduces. Mean over eight seeds. For reference, on the easier row-level split the raw mean calibration error is 0.103 and conformal reduces it to 0.013.}
\label{tab:calib}
\begin{tabular}{lcccc}
\toprule
Method & global & low $\log D$ & mid $\log D$ & high $\log D$ \\
\midrule
Raw & 0.72 & -- & -- & -- \\
Global conformal & 0.90 & 0.87 & 0.92 & 0.87 \\
Mondrian conformal & 0.87 & 0.96 & 0.82 & 0.88 \\
\bottomrule
\end{tabular}
\end{table}

\begin{table}[htbp]
\centering
\caption{Model and uncertainty-estimator comparison on a ligand-disjoint split (train/calibration/test $=55/20/25$ of ligands). Uncertainty-ranking quality is the Spearman correlation between predicted $\sigma$ and realised absolute error; higher is better. Coverage is the empirical coverage of split-conformal 80\% intervals. Mean over four seeds.}
\label{tab:uq}
\begin{tabular}{lcccc}
\toprule
Model & $R^2$ & conformal & rank quality & enrich. \\
 & & 80\% cov. & $\rho(|e|,\sigma)$ & factor \\
\midrule
Random forest & 0.15 & 0.91 & 0.22 & 7.5 \\
Quantile RF & 0.14 & 0.91 & 0.22 & 7.5 \\
Gaussian process & -0.03 & 0.95 & 0.23 & 3.7 \\
Gradient boosting & 0.22 & 0.91 & 0.28 & 7.5 \\
MC-dropout BNN & 0.12 & 0.94 & 0.06 & 6.6 \\
\bottomrule
\end{tabular}
\end{table}


\section{Reproduction}
All code is released at
\url{https://github.com/krishnatheaverage/logd-active-learning-benchmark} and archived
on Zenodo. The original benchmark is reproduced by \texttt{src/}: \texttt{real\_data.py}
(baseline), \texttt{real\_al.py} ($f$-element AL), \texttt{lipo\_al.py} (Lipophilicity
AL), \texttt{al\_calibration.py} (calibration), \texttt{prospective.py} (prospective
validation). The revision experiments are reproduced by \texttt{revision/}:
\texttt{data\_ext.py} (recovers ligand/source/lanthanide/condition group structure),
\texttt{e1\_splits.py} (leakage-controlled splits), \texttt{e2\_ligand\_batch.py}
(ligand-batch acquisition), \texttt{e3\_recall\_metrics.py} (recall definitions and
discovery metrics), \texttt{e4\_calibration.py} (subgroup calibration and Mondrian),
\texttt{e5\_uq\_baselines.py} (model/UQ comparison), with \texttt{make\_figures.py} and
\texttt{make\_tables.py} regenerating every figure and table. The $f$-element source
data are from Liu et al.\ (cited); the Lipophilicity set is from MoleculeNet.

\end{document}
